\documentclass[12pt, a4]{article}

\usepackage{amsmath}
\usepackage{german}
\usepackage{graphicx}
\DeclareMathOperator\sgn{sgn}

%\graphicspath{{/home/ritz/Teaching/fcn/fig/}}

%----------------- layout ----------------------------------------

\oddsidemargin=0in
\textwidth=6in
\topmargin=-0.7in
\textheight=9.5in
\baselineskip=16pt
\font\mittel=cmr17
\font\roman=cmr12
\font\klein=cmr9
\font\mgross=cmr10
\newfont{\CAP}{cmcsc10 scaled 900}
\pagestyle{empty}
\parindent=0mm

%----------------- macros ----------------------------------------
\newcommand{\conc}[1]{[{\rm C}]_{\text{#1}}}
\newcommand{\units}[1]{\text{#1}}
\newcommand{\kalium}{$\text{K}^+$}
\newcommand{\natrium}{$\text{Na}^+$}
\newcommand{\chlor}{$\text{Cl}^-$}
\newcommand{\wasser}{$\text{H}_2\text{O}$}
%-----------------------------------------------------------------


\begin{document}

\parbox{2cm}{
\includegraphics[width=1.8cm]{bccnlogo.pdf}
}
\parbox{11cm}{
\begin{center}
\large HUMBOLDT-UNIVERSIT\"AT \hskip 0.1 cm ZU \hskip 0.1 cm BERLIN
\vskip 0.1 true cm
\mgross BERNSTEIN CENTER FOR COMPUTATIONAL NEUROSCIENCE
\end{center}
}
\parbox{2cm}
{
\hfill
\includegraphics[width=1.8cm]{hublogo.pdf}
}

\vskip 0.6 true cm

\leftline{\CAP Humboldt-Universit\"at zu Berlin 
\hfill Phone: 030/2093-9110}
\leftline{\CAP Philippstr. 13 House 6
              \hfill Fax: 030/2093-6771}
\leftline{\CAP \hfill webpage: http://www.bccn-berlin.de/}


\vskip 0.8 true cm
\centerline{\bf Models of Neural Systems, WS 2010/11}
\centerline{\bf Computer Practical 5}
\centerline{Solutions to hand in on: November, 29th, 2010}

\vskip 0.8 true cm
{\bf Ordinary Differential Equations (ODEs)}

\medskip
Differential equations describe the evolution of systems in continous
time and are widely used in science and engineering. Here, we will
focus on the first order ordinary differential equations:
\begin{equation}
    \dot{x}(t) = f\left(x(t),t\right),
    \label{eq:ODE}
\end{equation} 

where $\dot{x}(t)\equiv\frac{\mathrm{d}x}{\mathrm{d}t}$ is the first derivative of
$x(t)$ with respect to time $t$. The solution to this equation
is the function $x(t)$ which satisfies the above relation.

\vskip 0.4 true cm
{\bf Exercises}
\begin{enumerate}

\item \textbf{Numerical solutions to ODEs.} The simplest numerical method for 
    solving the equation (\ref{eq:ODE}) is by discretization. Let's assume 
    that the function $x(t)$ is constant
    during the short time interval $\Delta t$. We can then rewrite the equation
    using finite steps:

    \begin{equation}
        \tilde{x}_{i+1} =\tilde{x}_i + f(\tilde{x}_i, i\Delta t)\Delta t,
        \quad \text{for } i=0,1,\dots,N.
        \label{eq:Euler}
    \end{equation} 

    In order to find the solution, we start with the value $x_0=x(0)$,
    called the initial condition, and then proceed by adding
    small increments to the function according to equation
    (\ref{eq:Euler}). This algorithm is called Euler method.

    \begin{enumerate}
        \item Define a Python function that calculates and returns a
        numerical solution to any first-order ODE. The function should
        take as an argument the function describing the right-hand side
	$f(x,t)$ of the ODE, an initial condition $x_0$, 
	an integration step $\Delta t$, and a stop time 
	$T_\mathrm{max}=N\Delta t$.
     
        \item Population growth can be modelled with the so-called logistic
        model: 
	\begin{equation} 
	    \frac{\mathrm{d}x}{\mathrm{d}t}=rx(1-x), \label{eq:logistic}
        \end{equation} where $r$ is a growth parameter. Use the Euler
        method to solve the logistic equation with different
        initial conditions and values of the parameter $r$.

        \item Plot the numerical solutions to the logistic model with
        different integration steps $\Delta t$ and compare them to the
        analytical solution, which is the sigma function:
        \begin{equation}
            x(t) = \left(1 + \left( \frac{1}{x_0} - 1 \right)
                \mathrm{exp}(-rt)\right)^{-1}.
        \end{equation}

        \item Calculate the mean square error between the
        numerical solutions $\tilde{x}_i$ and the analytical solutions $x(t)$
        as a function of $\Delta t$:
     
        \begin{equation}
            \mathrm{MSE}=\frac{1}{N}
            \sum_{i=0}^N (\tilde{x}_i-x(i\Delta t))^2.  \label{eq:MSE}
        \end{equation}
     
        Plot and explain the accuracy changes with the integration step
        $\Delta t$.
    \end{enumerate}


    \item \textbf{Passive membrane.}\label{ex:membrane} The following 
    equation describes the development of the membrane potential $V(t)$
    of a passive membrane, that is subject to the current injection $I(t)$:

    \begin{equation}
        \tau_m \frac{\mathrm{d}V(t)}{\mathrm{d}t}=-V(t)+E_m+R_mI(t).
        \label{eq:membrane}
    \end{equation} 

    At time $t=0$ the membrane voltage is at the resting potential $V(t=0)=E_m$.

    \begin{enumerate}
        \item Build a model cell from equation (\ref{eq:membrane}). Use the 
	    membrane resistance $R_m=10^7\,\mathrm{ \Omega}$, 
	    a membrane time constant of $\tau_m=10\mathrm{ms}$,
	    $I(t)=I_0=10^{-5}\,\mathrm{ mA}$,
            and $E_m=-80\,\mathrm{mV}$. Use the Euler
            method to find $V(t)$.
        \item Consider a time-dependent current injection, with a constant
            $10^{-5}\mathrm{mA}$ current between 10 and 100ms and zero current
            elsewise.
%         \begin{equation*}
%             I(t)= 
%             \begin{cases} 10^{-5}\,\mathrm{mA} & \text{if}\; 10\,\mathrm{ms}
%             \leq t\leq100\,\mathrm{ms},
%             \\
%             0 &\text{otherwise.}
%             \end{cases}
%         \end{equation*}
        Find the numerical solution for $V(t)$. Repeat the calculations for
        $\tau_m=1\mathrm{ms},100\mathrm{ms},1\mathrm{s}$.
    \end{enumerate}

    \item \textbf{Integrate-and-fire neuron.} Linear models are unable to 
	produce realistic action potentials, but
        spikes can be simulated with a simple reset mechanism.
        \begin{enumerate}

            \item Modify the Euler method so that every time at which the 
	    membrane potential reaches the threshold 
	    $V_\mathrm{th}=-54\,\mathrm{mV}$ the potential is set to the value
            $V_\mathrm{reset}=E_m$. The times of reset correspond to the spike 
	    times of the model.

            \item Use the integration with reset and simulate the neuron from
            Exercise \ref{ex:membrane} with a constant current intensity
            $I_0$. Plot the voltage traces for 500 ms simulation time.

            \item Calculate and plot the mean firing rate as a function of
            input intensity. How do you interpret the results?
        \end{enumerate}

    \item {\bf Report.} Submit the plain source code files for all exercises
        and a report containing the source code, all plots and answers to the
        questions asked.
%     and a report. The report should be in PDF format and include source
%     code to all exercises, answers to the
%     questions in exercises?, and plots from ? with
%     caption and labels.
\end{enumerate}

\vfill
\centerline{\CAP Contact}
\CAP

\begin{tabular}{lll}
Urs Bergmann & Phone: 2093-8924 & Email:
urs.bergmann(at)biologie.hu-berlin.de \\
Matthias Guggenmos & & Email: matthias.guggenmos(at)bccn-berlin.de \\
Richard Kempter \hfill & Phone: 2093-8925 \hfill & Email:
r.kempter(at)biologie.hu-berlin.de \\
Robert Schmidt & Phone: 2093-8926 & Email: r.schmidt(at)biologie.hu-berlin.de
\\
Bartosz Telenczuk & Phone: 2093-8838 & Email:
b.telenczuk(at)biologie.hu-berlin.de \\
\end{tabular}

\end{document}
 

